\chapter{Réalisation et Implémentation}
\label{chap:realisation}

\section{Introduction}

Ce chapitre décrit la réalisation concrète de l'application de gestion et de planification
des projets BTP développée pour CIVCO BTP, à partir de la conception présentée au
chapitre précédent. Nous présentons d'abord l'environnement technique et les outils
utilisés, puis détaillons l'implémentation du backend Laravel, du frontend React et
des principaux modules métier. Le chapitre se conclut par les tests réalisés et les
résultats de validation.

% ─────────────────────────────────────────────────────────────────────
\section{Environnement technique}

\subsection{Stack technologique}

\begin{description}
  \item[Serveur applicatif] PHP 8.3, Laravel 13.8, Laravel Sanctum 4
  (authentification SPA), Eloquent ORM, migrations et seeders.
  \item[Interface utilisateur] React 19.2, Vite 8, React Router 7, Tailwind CSS 4, Axios,
  Recharts (graphiques), Leaflet (cartographie).
  \item[Base de données] MySQL / PostgreSQL (modèle relationnel,
  schéma versionné par migrations Laravel).
  \item[Outils complémentaires] Git / GitHub (versionnement), Visual Studio Code (IDE),
  PHPUnit 12 (tests backend), ESLint (qualité frontend).
\end{description}

Le Tableau~\ref{tab:stack_versions} récapitule les versions principales des composants
utilisés en développement.

\begin{table}[H]
\caption{Versions des composants de la stack technique}
\label{tab:stack_versions}
\centering
\small
\begin{tabular}{|l|l|p{6.5cm}|}
\hline
\textbf{Composant} & \textbf{Version} & \textbf{Rôle} \\
\hline
PHP & 8.3+ & Langage backend, exécution Laravel \\
\hline
Laravel & 13.8 & Framework API, ORM, authentification \\
\hline
Laravel Sanctum & 4.x & Sessions SPA sécurisées \\
\hline
React & 19.2 & Interface utilisateur composants \\
\hline
Vite & 8.0 & Build tool et serveur de développement \\
\hline
Tailwind CSS & 4.3 & Styles utilitaires \\
\hline
MySQL / PostgreSQL & 8.x / 16 & Persistance relationnelle \\
\hline
PHPUnit & 12.5 & Tests automatisés backend \\
\hline
\end{tabular}
\end{table}

\subsection{Outils de développement}

\begin{description}
  \item[Gestion de projet] Trello (tableau Kanban) pour planifier les modules,
  suivre l'avancement des tâches et préparer les démonstrations de suivi.
  \item[Contrôle de version] Git et GitHub pour le suivi des évolutions, branches de
  fonctionnalités et revues de code.
  \item[IDE] Visual Studio Code avec extensions PHP Intelephense, ESLint et Prettier.
  \item[Tests] PHPUnit pour les tests d'intégration API ; tests manuels sur navigateur
  (Chrome, Edge) pour les parcours utilisateurs.
  \item[Qualité de code] Laravel Pint (formatage PHP), ESLint (linting JavaScript/React).
  \item[Documentation API] Routes Laravel exposées sous \texttt{/api/v1/}, testables via
  outils HTTP (Postman, curl) ou le navigateur pour les endpoints GET.
\end{description}

\subsection{Architecture de déploiement}

\subsubsection{Environnement de développement}

En développement, l'application repose sur deux processus parallèles, orchestrés depuis
deux terminaux :

\begin{enumerate}
  \item \textbf{Backend Laravel} : \texttt{php artisan serve --host=127.0.0.1 --port=8000}
  expose l'API sur le port 8000.
  \item \textbf{Frontend Vite} : \texttt{npm run dev} lance le serveur de développement
  React sur le port 5173, avec rechargement à chaud (\textit{hot-reload}).
\end{enumerate}

Le fichier \texttt{vite.config.js} configure un proxy : les requêtes vers \texttt{/api},
\texttt{/sanctum} et \texttt{/storage} sont redirigées vers Laravel (port 8000), ce qui
évite les problèmes de CORS en développement et simule un déploiement unifié.

\subsubsection{Schéma de communication}

\begin{itemize}
  \item Le navigateur charge l'application React depuis \texttt{http://127.0.0.1:5173}.
  \item Les appels API transitent par le proxy Vite vers \texttt{http://127.0.0.1:8000/api/v1/}.
  \item Sanctum émet un cookie de session après authentification ; Axios est configuré
  avec \texttt{withCredentials: true}.
  \item Les fichiers uploadés (logos, documents, signatures) sont stockés dans
  \texttt{storage/app/public} et servis via le lien symbolique \texttt{public/storage}.
\end{itemize}

\subsubsection{Étude d'hébergement : OVHcloud vs o2switch}

En vue d'une mise en production chez CIVCO BTP, une \textbf{étude comparative
d'hébergement} a été menée entre deux prestataires français largement utilisés pour
les applications web PHP : \textbf{OVHcloud} et \textbf{o2switch}. L'objectif était
d'identifier la solution la mieux adaptée au profil de l'entreprise : application
métier interne, équipe réduite, budget maîtrisé et besoin de simplicité
d'administration.

\medskip

Les critères d'évaluation retenus sont les suivants :

\begin{itemize}
  \item \textbf{Compatibilité Laravel / PHP 8.3} : exécution de l'API backend et des
  commandes \texttt{artisan} (migrations, files d'attente, planificateur).
  \item \textbf{Base de données} : support MySQL ou PostgreSQL, sauvegardes et quotas.
  \item \textbf{Déploiement du frontend React} : hébergement des fichiers statiques
  produits par \texttt{npm run build}.
  \item \textbf{Sécurité} : certificat SSL, pare-feu, mises à jour système.
  \item \textbf{Coût et prévisibilité} : tarification mensuelle, absence de surprises
  sur la facturation.
  \item \textbf{Support et localisation} : assistance en français, datacenters en
  France, conformité pour des données clients BTP.
  \item \textbf{Facilité d'exploitation} : courbe d'apprentissage pour une PME sans
  équipe DevOps dédiée.
\end{itemize}

\begin{table}[H]
\caption{Comparatif d'hébergement : OVHcloud vs o2switch}
\label{tab:hebergement}
\centering
\small
\begin{tabular}{|p{3.2cm}|p{5.2cm}|p{5.2cm}|}
\hline
\textbf{Critère} & \textbf{OVHcloud} & \textbf{o2switch} \\
\hline
Type d'offre & VPS, cloud public, serveurs dédiés (IaaS) & Hébergement mutualisé
tout-en-un \\
\hline
Laravel / PHP & Compatible (configuration manuelle : Nginx, PHP-FPM, Composer) &
Compatible nativement (PHP 8.x, SSH, Composer) \\
\hline
Base de données & MySQL/PostgreSQL à installer et administrer & MySQL inclus,
bases multiples \\
\hline
Frontend React (build) & Fichiers statiques sur VPS ou Object Storage & Hébergement
direct des assets \texttt{dist/} \\
\hline
SSL (HTTPS) & Let's Encrypt ou certificat payant (selon offre) & Let's Encrypt
inclus et automatisé \\
\hline
Sauvegardes & À configurer (snapshots, scripts) & Sauvegardes journalières incluses \\
\hline
Tarification & Variable (facturation à l'usage, snapshots, trafic) & Forfait fixe
illimité (offre unique) \\
\hline
Support & Tickets, documentation technique & Support français réactif, panel cPanel \\
\hline
Scalabilité & Élevée (montée en charge horizontale) & Modérée (mutualisé, adapté PME) \\
\hline
Complexité admin & Élevée (sysadmin requis) & Faible (panel graphique) \\
\hline
\end{tabular}
\end{table}

\medskip

\textbf{Analyse.} OVHcloud~\cite{ovhcloud} offre une grande flexibilité et une
scalabilité adaptée aux applications à fort trafic ou aux architectures distribuées
(conteneurs Docker, load balancing, bases managées). En revanche, le déploiement d'une
application Laravel + React sur un VPS OVHcloud suppose une configuration manuelle
du serveur web, de PHP-FPM, des certificats, des sauvegardes et de la supervision.
Cette approche est pertinente pour des équipes disposant de compétences
d'administration système, mais elle représente un surcoût en temps pour une PME comme
CIVCO BTP.

\medskip

o2switch~\cite{o2switch}, hébergeur français basé à Grabels, propose une offre
mutualisée orientée simplicité : PHP récent, bases MySQL, certificats SSL gratuits,
sauvegardes automatiques et support technique francophone. Pour une application métier
interne au volume d'utilisateurs modéré (équipes CIVCO BTP et clients portail), cette
formule couvre les besoins sans nécessiter la gestion d'un serveur virtuel. Le frontend
React est déployé sous forme de fichiers statiques (\texttt{npm run build}), tandis que
Laravel gère l'API et les sessions Sanctum sur le même domaine, simplifiant la
configuration des cookies et du CSRF.

\medskip

\textbf{Choix retenu.} Au terme de cette étude, \textbf{o2switch} a été retenu comme
plateforme d'hébergement cible pour CIVCO BTP. Ce choix se justifie par :

\begin{itemize}
  \item un \textbf{rapport qualité/prix} favorable pour une application interne ;
  \item une \textbf{mise en production plus rapide}, sans administration système
  avancée ;
  \item une \textbf{compatibilité native} avec la stack Laravel / MySQL ;
  \item un \textbf{support en français} et des datacenters localisés en France ;
  \item des \textbf{sauvegardes et SSL inclus}, répondant aux exigences de sécurité
  identifiées au Chapitre~3.
\end{itemize}

OVHcloud reste une alternative crédible si l'application devait évoluer vers une
architecture plus complexe (montée en charge importante, microservices, conteneurisation
Docker). Pour la première version de l'outil, la simplicité opérationnelle d'o2switch
a été privilégiée.

\subsubsection{Configuration : variables d'environnement}

La configuration sensible est centralisée dans le fichier \texttt{.env} du backend.
Le Tableau~\ref{tab:env_vars} liste les variables principales :

\begin{table}[H]
\caption{Variables d'environnement principales du fichier \texttt{.env}}
\label{tab:env_vars}
\centering
\small
\begin{tabular}{|l|p{7.5cm}|}
\hline
\textbf{Variable} & \textbf{Rôle} \\
\hline
\texttt{APP\_URL} & URL publique de l'application (ex. \texttt{http://127.0.0.1:8000}) \\
\hline
\texttt{DB\_CONNECTION / DB\_DATABASE} & Connexion SGBD (MySQL ou PostgreSQL) \\
\hline
\texttt{SESSION\_DRIVER} & Pilote de session (fichier ou base de données) \\
\hline
\texttt{SANCTUM\_STATEFUL\_DOMAINS} & Domaines autorisés pour les cookies Sanctum \\
\hline
\texttt{FRONTEND\_URL} & URL du frontend (pour les redirections et CORS) \\
\hline
\texttt{MAIL\_*} & Configuration SMTP (notifications, invitations portail) \\
\hline
\end{tabular}
\end{table}

\subsubsection{Initialisation de la base de données}

Le démarrage d'un environnement neuf suit la séquence standard Laravel :

\begin{enumerate}
  \item \texttt{composer install} : installation des dépendances PHP.
  \item \texttt{cp .env.example .env} puis configuration de la base de données.
  \item \texttt{php artisan key:generate} : génération de la clé d'application.
  \item \texttt{php artisan migrate --seed} : création du schéma et données de démonstration.
  \item \texttt{php artisan storage:link} : lien symbolique pour les fichiers publics.
\end{enumerate}

Les seeders (\texttt{DatabaseSeeder}, \texttt{DemoSeeder}, \texttt{RoleSeeder}) peuplent
la base avec des utilisateurs de test, des rôles, des permissions, des clients et des
projets représentatifs du secteur BTP au Maroc.

% ─────────────────────────────────────────────────────────────────────
\section{Implémentation du backend}

\subsection{Structure du projet}

Le backend suit l'architecture modulaire Laravel, organisée par responsabilité :

\begin{verbatim}
backend/
+-- app/
|   +-- Enums/              # Statuts (Project, Quote, Invoice, Client...)
|   +-- Http/
|   |   +-- Controllers/Api/V1/   # ~45 controleurs REST
|   |   +-- Middleware/           # Auth, permissions, contexte
|   |   +-- Requests/             # Validation par domaine
|   |   +-- Resources/            # Serialisation JSON
|   +-- Models/             # Entites Eloquent (~35 modeles)
|   +-- Services/           # Logique metier (QuoteService, etc.)
+-- database/
|   +-- migrations/         # ~70 migrations versionnees
|   +-- seeders/            # Donnees de demo et referentiels
+-- routes/api.php          # Definition des routes /api/v1/
\-- tests/Feature/          # Tests PHPUnit d'integration
\end{verbatim}

Cette organisation sépare clairement le transport HTTP (contrôleurs), la validation
(Requests), la transformation des réponses (Resources) et la logique métier (Services),
conformément au patron décrit au Chapitre~3.

\subsection{API REST : endpoints principaux}

L'API est regroupée sous le préfixe \texttt{/api/v1/}. Les endpoints principaux sont
organisés par domaine fonctionnel :

\begin{description}
  \item[Authentification] \texttt{POST /login}, \texttt{GET /me}, \texttt{POST /logout},
  \texttt{PATCH /me} (profil utilisateur). Gestion de session via Sanctum.
  \item[Projets] \texttt{GET/POST /projects}, \texttt{GET/PATCH/DELETE /projects/\{id\}},
  gestion des phases, équipe, dépenses, médias et commentaires.
  \item[Clients] \texttt{GET/POST /clients}, archivage, contacts, provisionnement
  portail (\texttt{POST /clients/\{id\}/portal-account}).
  \item[Devis] \texttt{GET/POST /quotes}, lignes de détail, conversion en facture,
  suivi des statuts et impressions.
  \item[Factures] \texttt{GET/POST /invoices}, paiements, état de facturation.
  \item[Bons de livraison] \texttt{GET/POST /delivery-forms}, lignes et statuts.
  \item[Contrats] \texttt{GET/POST /contracts}, compilation, signature, avenants
  (\texttt{/contract-amendments}).
  \item[Tâches] \texttt{GET/POST /tasks}, assignation, filtrage par projet et utilisateur.
  \item[Équipe] \texttt{GET /team/members}, gestion des membres et rôles.
  \item[Configuration] \texttt{/roles}, \texttt{/document-templates},
  \texttt{/tenant/document-controls}, thème et badges.
  \item[Tableau de bord] \texttt{GET /dashboard/summary}, indicateurs agrégés.
  \item[Recherche] \texttt{GET /search}, recherche globale multi-entités.
  \item[Portail client] \texttt{/client-portal/*} : projets, devis, contrats,
  messagerie, acceptation et signature.
  \item[Historique] \texttt{GET /activity-logs}, journal des actions.
\end{description}

Chaque route métier sensible est protégée par le middleware \texttt{permission:*},
qui vérifie que l'utilisateur authentifié dispose du droit requis avant d'exécuter
l'action.

\subsection{Services métier}

La logique complexe est extraite des contrôleurs vers des classes Service :

\begin{itemize}
  \item \textbf{QuoteService / InvoiceService} : calcul des totaux, numérotation,
  conversion devis $\rightarrow$ facture, gestion des statuts.
  \item \textbf{DeliveryFormService} : production des bons de livraison et liaison
  aux projets.
  \item \textbf{DocumentService / DocumentTemplateCompilationService} : génération
  et compilation des documents à partir de modèles.
  \item \textbf{DashboardService} : agrégation des indicateurs (projets actifs,
  tâches en retard, chiffre d'affaires).
  \item \textbf{AuthContextService} : résolution du contexte utilisateur (rôles,
  permissions, société courante).
  \item \textbf{ActivityLogService} : journalisation des actions sensibles.
  \item \textbf{PrintTrackingService} : suivi des impressions officielles et quotas.
  \item \textbf{SiteGeocodingService} : géolocalisation des chantiers (Nominatim +
  coordonnées de repli pour les villes marocaines).
\end{itemize}

\subsection{Authentification et contrôle d'accès}

L'authentification repose sur \textbf{Laravel Sanctum} en mode SPA :

\begin{enumerate}
  \item Le frontend obtient un cookie CSRF via \texttt{GET /sanctum/csrf-cookie}.
  \item L'utilisateur envoie ses identifiants à \texttt{POST /api/v1/login}.
  \item Laravel valide, crée une session et retourne le profil (\texttt{UserResource})
  avec rôles et permissions.
  \item Chaque requête ultérieure inclut le cookie de session ; le middleware
  \texttt{auth:sanctum} protège les routes.
  \item Le middleware \texttt{permission:slug} vérifie les droits granulaires
  (\texttt{project.view}, \texttt{quote.manage}, etc.).
\end{enumerate}

Le service \texttt{PostLoginRedirectService} oriente l'utilisateur vers la page
d'accueil adaptée à son profil (tableau de bord admin, liste des projets pour un
chef de projet, factures pour un comptable).

\subsection{Migrations et modèle de données}

Le schéma de base compte environ 70 migrations, couvrant l'ensemble des entités métier.
Les migrations sont idempotentes et versionnées : chaque modification de schéma
produit un nouveau fichier daté, exécutable de façon reproductible sur tout
environnement.

\medskip

Les relations Eloquent principales ont été implémentées comme suit :
\begin{itemize}
  \item \texttt{Client} \texttt{hasMany} \texttt{Project}
  \item \texttt{Project} \texttt{belongsTo} \texttt{Client}, \texttt{hasMany} \texttt{ProjectPhase},
  \texttt{Task}, \texttt{Quote}, \texttt{Contract}
  \item \texttt{Quote} \texttt{hasMany} \texttt{QuoteLine}, \texttt{hasOne} \texttt{Invoice}
  \item \texttt{User} \texttt{belongsToMany} \texttt{Project} (équipe), \texttt{belongsTo} \texttt{Role}
\end{itemize}

Les \textbf{API Resources} (\texttt{ProjectResource}, \texttt{QuoteResource}, etc.)
contrôlent les champs exposés au frontend et permettent d'inclure conditionnellement
les relations (\texttt{client}, \texttt{lines}, \texttt{teamMembers}).

\subsection{Middleware et sécurité des routes}

Le pipeline de middleware Laravel applique une chaîne de vérifications sur chaque
requête API protégée :

\begin{enumerate}
  \item \textbf{auth:sanctum} : vérifie que l'utilisateur est authentifié via session.
  \item \textbf{EnsureCompanyContext} : résout la société courante à partir de
  l'en-tête \texttt{X-Company-Id} ou du contexte utilisateur.
  \item \textbf{permission:slug} : contrôle granulaire par action métier
  (\texttt{client.view}, \texttt{invoice.manage}, etc.).
  \item \textbf{ValidateCsrfToken} : protection contre les requêtes forgées sur
  les verbes mutants (\texttt{POST}, \texttt{PUT}, \texttt{PATCH}, \texttt{DELETE}).
\end{enumerate}

Cette approche en profondeur (\textit{defense in depth}) garantit qu'aucune action
métier ne peut être exécutée sans authentification et sans autorisation explicite,
conformément aux bonnes pratiques de sécurité web~\cite{owasp}.

\subsection{Exemple : création d'un devis}

Le flux de création d'un devis illustre l'interaction entre les couches :

\begin{enumerate}
  \item Le frontend envoie \texttt{POST /api/v1/quotes} avec le corps JSON
  (client, projet, lignes, montants).
  \item \texttt{StoreQuoteRequest} valide les champs (références obligatoires,
  montants positifs, dates cohérentes).
  \item \texttt{QuoteController@store} délègue à \texttt{QuoteService::create()}.
  \item Le service crée l'enregistrement \texttt{Quote} et ses \texttt{QuoteLine}
  associées dans une transaction base de données.
  \item Les totaux HT, TVA et TTC sont calculés automatiquement.
  \item \texttt{QuoteResource} sérialise la réponse JSON pour le frontend.
  \item \texttt{ActivityLogService} journalise l'action si configuré.
\end{enumerate}

Ce patron (Request $\rightarrow$ Controller $\rightarrow$ Service $\rightarrow$
Resource) est reproduit sur l'ensemble des modules métier.

\subsection{Modules complémentaires}

En complément des modules principaux, plusieurs fonctionnalités transverses ont été
implémentées :

\medskip

\textbf{Notifications :} le modèle \texttt{Notification} et le contrôleur associé
permettent d'informer les utilisateurs des événements importants (nouvelle tâche
assignée, devis accepté, message reçu). L'interface affiche un badge de notification
non lue dans l'en-tête.

\medskip

\textbf{Messagerie :} le module \texttt{PortalMessage} gère les échanges entre
l'équipe interne et les clients externes, scopés par projet. Le service
\texttt{PortalMessageService} filtre les conversations selon les droits de
l'utilisateur (\texttt{can\_chat\_with\_client} sur les membres d'équipe).

\medskip

\textbf{Recherche globale :} le composant \texttt{GlobalSearch} et l'endpoint
\texttt{GET /search} permettent de retrouver rapidement un client, un projet, un
devis ou une facture par mot-clé, améliorant la productivité au quotidien.

\medskip

\textbf{Contrôle d'impression :} \texttt{PrintTrackingService} gère les quotas
d'impressions officielles par type de document et par en-tête (avec ou sans papier
à en-tête). Au-delà du quota, les impressions sont marquées « COPIE - NON OFFICIEL ».

\medskip

\textbf{Journal d'activité :} chaque action sensible (création de projet, validation
de devis, archivage client) est enregistrée par \texttt{ActivityLogService} avec
l'identité de l'utilisateur, le type d'action et l'horodatage. Le module
\texttt{features/history/} expose ces entrées sous forme de fil d'audit consultable
par l'administrateur.

\medskip

\textbf{Géocodage des chantiers :} lors de la saisie d'une adresse de chantier,
\texttt{SiteGeocodingService} tente de résoudre automatiquement les coordonnées
GPS via l'API Nominatim (OpenStreetMap). En cas d'échec, un référentiel de
coordonnées des principales villes marocaines (\texttt{MoroccoCityCoordinates})
assure un positionnement approximatif sur la carte Leaflet.

\medskip

\textbf{Modèles de documents :} le module \texttt{DocumentTemplate} permet de
définir des gabarits de contrats et de documents commerciaux avec des champs
dynamiques (placeholders). Le service \texttt{DocumentTemplateCompilationService}
injecte les données du projet et du client lors de la génération du document final.

\medskip

\textbf{Gestion d'équipe :} \texttt{TeamMemberService} et \texttt{TeamQueryService}
centralisent l'affectation des collaborateurs aux projets, la gestion des rôles
sur chantier et le contrôle de l'autorisation de messagerie client
(\texttt{can\_chat\_with\_client}).

% ─────────────────────────────────────────────────────────────────────
\section{Implémentation du frontend}

\subsection{Architecture React}

Le frontend est une SPA React 19 organisée par domaines métier (\textit{features}) :

\begin{verbatim}
frontend/src/
+-- api/                    # Modules Axios par ressource
|   +-- client.js           # Instance Axios + intercepteurs
|   +-- projects.js, clients.js, quotes.js, invoices.js...
+-- components/             # Composants transverses (Sidebar, Modal...)
+-- context/
|   +-- AuthContext.jsx     # Session, roles, permissions
|   +-- ThemeContext.jsx    # Theme sombre, couleurs societe
|   +-- LanguageContext.jsx # i18n FR / EN
+-- features/
|   +-- dashboard/          # Tableau de bord et widgets
|   +-- projects/           # Liste, detail, phases, equipe
|   +-- clients/            # Gestion clients et portail
|   +-- quotes/             # Devis
|   +-- invoices/           # Factures
|   +-- delivery-forms/     # Bons de livraison
|   +-- tasks/              # Taches
|   +-- configuration/      # Parametres, roles, modeles
|   +-- profile/            # Profil utilisateur
+-- routes/
|   +-- AppRoutes.jsx       # Definition des routes
|   +-- routePermissions.js # Controle d'acces par route
\-- i18n/locales/           # Traductions fr.js, en.js
\end{verbatim}

\subsection{Client HTTP et authentification}

Le module \texttt{api/client.js} centralise la configuration Axios :

\begin{itemize}
  \item \texttt{baseURL} pointant vers \texttt{/api/v1} (proxifié par Vite).
  \item \texttt{withCredentials: true} pour transmettre les cookies Sanctum.
  \item Intercepteur de requête : récupération du jeton CSRF depuis le cookie
  \texttt{XSRF-TOKEN} et injection dans l'en-tête \texttt{X-XSRF-TOKEN}.
  \item Intercepteur de réponse : redirection vers \texttt{/login} en cas de
  \texttt{401 Unauthorized}.
  \item En-tête \texttt{X-Company-Id} pour le contexte société lorsque requis.
\end{itemize}

Le \texttt{AuthContext} expose l'état global d'authentification : utilisateur courant,
rôles, permissions, fonctions \texttt{login} / \texttt{logout}, et helpers de
vérification des droits consommés par les composants \texttt{PermissionGate} et
\texttt{AdminRoute}.

\subsection{Routage et contrôle d'accès}

\texttt{AppRoutes.jsx} définit l'arborescence de navigation :

\begin{itemize}
  \item Routes publiques : \texttt{/login}.
  \item Routes authentifiées : \texttt{/} (dashboard), \texttt{/projects},
  \texttt{/clients}, \texttt{/quotes}, \texttt{/invoices}, \texttt{/tasks},
  \texttt{/map}, \texttt{/configuration}, \texttt{/history}.
  \item Portail client : \texttt{/portal/*} (interface dédiée aux clients externes).
  \item Protection : composants \texttt{ProtectedRoute}, \texttt{AdminRoute} et
  fonction \texttt{canAccessRoute()} qui croise le chemin demandé avec les
  permissions de l'utilisateur.
\end{itemize}

\subsection{Composants transverses}

Plusieurs composants partagés structurent l'expérience utilisateur :

\begin{itemize}
  \item \textbf{Sidebar} : navigation principale, filtrée selon les permissions.
  \item \textbf{AppLayout} : enveloppe commune (sidebar + en-tête + contenu).
  \item \textbf{GlobalSearch} : recherche unifiée clients, projets, documents.
  \item \textbf{CommercialPrintSheet} : génération et impression des documents
  commerciaux (devis, factures, BL) avec suivi des impressions.
  \item \textbf{PrintOptionsModal} : choix d'impression (avec/sans en-tête,
  copie officielle).
\end{itemize}

\subsection{Internationalisation}

L'application supporte le français (par défaut) et l'anglais via \texttt{LanguageContext}
et les fichiers \texttt{i18n/locales/fr.js} et \texttt{en.js}. Les libellés de
l'interface, les messages d'erreur et les intitulés de menus sont externalisés,
facilitant l'ajout de nouvelles langues.

\subsection{Détail des modules frontend}

Chaque module \textit{feature} suit une structure cohérente : une page liste
(\texttt{*Page.jsx}), une ou plusieurs pages détail (\texttt{*DetailPage.jsx}),
des composants locaux (\texttt{components/}) et un module API dédié
(\texttt{api/*.js}).

\medskip

\textbf{Module Projets} (\texttt{features/projects/}) : \texttt{ProjectsPage} affiche
la liste filtrable des chantiers ; \texttt{ProjectDetailPage} regroupe les onglets
phases, tâches, équipe, documents et contrats. Le composant \texttt{ProjectAmendmentsTab}
gère les avenants contractuels.

\medskip

\textbf{Module Clients} (\texttt{features/clients/}) : \texttt{ClientsPage} avec
recherche et filtres ; \texttt{NewClientModal} pour la création ; \texttt{ClientPortalPanel}
pour la gestion de l'accès portail. Le composant \texttt{ConfirmArchiveModal} sécurise
l'archivage.

\medskip

\textbf{Module Documents commerciaux} : \texttt{QuotesPage}, \texttt{QuoteDetailPage},
\texttt{InvoicesPage}, \texttt{DeliveryFormDetailPage}. Les composants d'impression
(\texttt{CommercialPrintSheet}, \texttt{DeliveryFormPrintSheet}) génèrent les
versions PDF imprimables.

\medskip

\textbf{Module Configuration} (\texttt{features/configuration/}) : panneaux dédiés
pour les rôles (\texttt{RolesSettingsPanel}), les modèles de contrats
(\texttt{ContractTemplateSettingsPanel}), les contrôles d'impression
(\texttt{DocumentControlsSettingsPanel}) et les modèles de documents
(\texttt{DocumentTemplatesSettingsPanel}).

\subsection{Gestion d'état et hooks}

L'état applicatif repose principalement sur React Context :

\begin{itemize}
  \item \textbf{AuthContext} : session utilisateur, permissions, fonctions login/logout.
  \item \textbf{ThemeContext} : couleurs de la société injectées en variables CSS.
  \item \textbf{LanguageContext} : bascule FR / EN de l'interface.
  \item \textbf{usePolicyPrint / useProjectTasks} : hooks métier encapsulant la
  logique d'impression et de chargement des tâches projet.
\end{itemize}

React Query (\texttt{@tanstack/react-query}) est installé et prêt pour le cache
des requêtes API ; son adoption progressive permettra d'optimiser les rechargements
de listes et fiches détail.

% ─────────────────────────────────────────────────────────────────────
\section{Déroulement du développement et livrables}

\subsection{Phases de réalisation}

Le développement a suivi la planification définie au Chapitre~1, avec une
progression par modules fonctionnels. Chaque phase a produit un incrément
démontrable lors des réunions de suivi hebdomadaires avec l'encadrant professionnel.

\begin{table}[H]
\caption{Phases de développement et modules livrés}
\label{tab:livrables}
\centering
\small
\begin{tabular}{|l|p{5.5cm}|p{5.5cm}|}
\hline
\textbf{Phase} & \textbf{Modules livrés} & \textbf{Validation} \\
\hline
Fondations & Auth Sanctum, RBAC, navigation React, layout & Connexion multi-profils \\
\hline
Cœur métier & Projets, clients, phases, équipe & CRUD complet, carte \\
\hline
Documents & Devis, factures, BL, impression & Cycle commercial bout en bout \\
\hline
Opérations & Tâches, contrats, avenants, dépenses & Suivi opérationnel \\
\hline
Collaboration & Portail client, messagerie, notifications & Parcours client externe \\
\hline
Gouvernance & Configuration, historique, recherche globale & Administration complète \\
\hline
\end{tabular}
\end{table}

\subsection{Synthèse quantitative des livrables}

Au terme du stage, l'application comprend les éléments quantitatifs suivants,
illustrant l'ampleur de la réalisation :

\begin{itemize}
  \item Plus de \textbf{70 migrations} de base de données versionnées.
  \item Environ \textbf{40 contrôleurs API} et \textbf{25 services métier}.
  \item Plus de \textbf{30 modules frontend} (\textit{features} et composants).
  \item Une suite de \textbf{7 classes de tests} PHPUnit couvrant les flux critiques.
  \item Plus de \textbf{200 routes API} protégées par middleware d'authentification
  et de permissions.
  \item Support \textbf{bilingue} (français / anglais) sur l'ensemble de l'interface.
\end{itemize}

\subsection{Bonnes pratiques adoptées}

Plusieurs bonnes pratiques de génie logiciel ont été appliquées tout au long du
développement :

\begin{itemize}
  \item \textbf{Séparation des responsabilités} : contrôleurs fins, logique métier
  dans les services, validation dans les Form Requests.
  \item \textbf{Conventions de nommage} : respect des standards Laravel (PascalCase
  pour les classes, snake\_case pour les colonnes) et React (PascalCase pour les
  composants).
  \item \textbf{Versionnement Git} : branches par fonctionnalité, commits atomiques,
  messages descriptifs.
  \item \textbf{Suivi Kanban (Trello)} : chaque module ou correction était modélisé
  sous forme de carte, déplacée au fil de l'avancement entre les colonnes du tableau.
  \item \textbf{Revue de code} : validation des choix techniques avec l'encadrant
  professionnel avant intégration des modules sensibles (facturation, permissions).
  \item \textbf{Documentation inline} : commentaires sur les règles métier non
  évidentes, docblocks sur les méthodes de service exposées à l'API.
\end{itemize}

% ─────────────────────────────────────────────────────────────────────
\section{Interfaces principales}

Cette section présente les écrans clés de l'application. Les captures d'écran seront
insérées ultérieurement ; les descriptions ci-dessous reflètent l'implémentation
actuelle.

\subsection{Page de connexion}

\begin{figure}[H]
  \centering
  \includegraphics[width=\textwidth]{login.png}
  \caption{Page de connexion : authentification par email et mot de passe}
  \label{fig:login}
\end{figure}

La page de connexion (Figure~\ref{fig:login}) permet à tout utilisateur interne de
CIVCO BTP de s'authentifier. Après validation, l'utilisateur est redirigé vers la
page d'accueil correspondant à son profil métier.

\medskip

L'écran affiche le logo de l'entreprise, un formulaire email/mot de passe et une
option « Se souvenir de moi ». Les erreurs d'authentification (identifiants
incorrects, compte désactivé) sont affichées de manière explicite. Le design sombre
et épuré a été choisi pour réduire la fatigue visuelle lors d'une utilisation
quotidienne prolongée.

\subsection{Tableau de bord}

\begin{figure}[H]
  \centering
  \includegraphics[width=\textwidth]{dashboard.png}
  \caption{Tableau de bord : vue synthétique de l'activité (projets, tâches, indicateurs)}
  \label{fig:dashboard}
\end{figure}

Le tableau de bord (Figure~\ref{fig:dashboard}) agrège les informations essentielles :
nombre de projets en cours, tâches prioritaires, graphiques d'avancement (Recharts)
et accès rapide aux modules. Les widgets affichés dépendent des permissions de
l'utilisateur connecté.

\medskip

Les widgets sont configurés dynamiquement : un administrateur voit l'ensemble des
indicateurs (projets actifs, factures en attente, tâches en retard), tandis qu'un
technicien ne voit que ses tâches assignées et les projets auxquels il participe.
Les graphiques Recharts illustrent l'évolution de l'avancement et la répartition
des projets par statut.

\subsection{Gestion des projets}

\begin{figure}[H]
  \centering
  \includegraphics[width=\textwidth]{projects.png}
  \caption{Module Projets : liste des chantiers et fiche détail (phases, équipe, documents)}
  \label{fig:projects}
\end{figure}

Le module Projets (Figure~\ref{fig:projects}) permet de créer, consulter et modifier
les chantiers. La fiche détail est organisée en onglets : informations générales,
phases, tâches, équipe, devis/factures liés, contrats, dépenses et médias.

\medskip

La liste des projets propose des filtres par statut, client et nature de chantier.
La fiche détail centralise toutes les informations d'un chantier : depuis un seul
écran, le chef de projet peut consulter l'avancement des phases, réaffecter un
collaborateur, joindre un document ou créer un devis lié au projet.

\subsection{Gestion des clients}

\begin{figure}[H]
  \centering
  \includegraphics[width=\textwidth]{clients.png}
  \caption{Module Clients : fiches clients, contacts et archivage}
  \label{fig:clients}
\end{figure}

Le module Clients (Figure~\ref{fig:clients}) centralise les maîtres d'ouvrage et
clients commerciaux. L'administrateur peut provisionner un accès portail pour un
client externe directement depuis la fiche client.

\medskip

Chaque fiche client regroupe les coordonnées, les contacts associés, les badges
métier (Gros Œuvre, Urgent, etc.), l'historique des projets et l'état du compte
portail. L'archivage est sécurisé par une modale de confirmation pour éviter les
suppressions accidentelles.

\subsection{Documents commerciaux}

\begin{figure}[H]
  \centering
  \includegraphics[width=\textwidth]{quotes.png}
  \caption{Module Devis : création avec lignes de détail et calcul automatique des totaux}
  \label{fig:quotes}
\end{figure}

\begin{figure}[H]
  \centering
  \includegraphics[width=\textwidth]{invoices.png}
  \caption{Module Factures : suivi des factures et état de paiement}
  \label{fig:invoices}
\end{figure}

Les modules Devis et Factures (Figures~\ref{fig:quotes} et~\ref{fig:invoices})
permettent la production des documents commerciaux, leur impression (avec suivi
des copies officielles) et la conversion devis $\rightarrow$ facture en un clic.

\medskip

Le formulaire de création de devis permet d'ajouter dynamiquement des lignes de
prestation (description, quantité, prix unitaire HT, taux de TVA). Les totaux sont
recalculés en temps réel. L'impression s'effectue via un composant React dédié
(\texttt{CommercialPrintSheet}) qui génère une version formatée prête à l'envoi au
client, avec suivi du quota d'impressions officielles.

\subsection{Planification des tâches}

\begin{figure}[H]
  \centering
  \includegraphics[width=\textwidth]{tasks.png}
  \caption{Module Tâches : planification, assignation et suivi par projet}
  \label{fig:tasks}
\end{figure}

Le module Tâches (Figure~\ref{fig:tasks}) offre une vue tableau des missions
opérationnelles, filtrable par statut, projet et assigné. Les chefs de projet
peuvent créer et assigner des tâches ; les utilisateurs internes consultent
leurs tâches propres.

\medskip

Le module propose une vue tableau avec tri par date d'échéance, filtre par statut
(à faire, en cours, terminé) et code couleur pour les tâches en retard. La création
d'une tâche associe obligatoirement un projet, une phase et un responsable, garantissant
la traçabilité de l'assignation.

\subsection{Carte des chantiers}

\begin{figure}[H]
  \centering
  \includegraphics[width=\textwidth]{map.png}
  \caption{Carte des projets : visualisation géographique des chantiers (Leaflet)}
  \label{fig:map}
\end{figure}

La carte (Figure~\ref{fig:map}) affiche les projets géolocalisés sur une carte
interactive Leaflet, utile pour visualiser la répartition géographique des
chantiers de CIVCO BTP.

\medskip

Chaque marqueur sur la carte correspond à un projet géolocalisé. Un clic affiche un
résumé (référence, client, statut, avancement) et un lien vers la fiche détail.
Cette visualisation est particulièrement utile pour la direction et les chefs de
projet qui gèrent plusieurs chantiers répartis sur différentes zones géographiques
du Grand Casablanca et au-delà.

\subsection{Portail client}

\begin{figure}[H]
  \centering
  \includegraphics[width=\textwidth]{portal.png}
  \caption{Portail client : consultation des projets, acceptation de devis et signature}
  \label{fig:portal}
\end{figure}

Le portail client (Figure~\ref{fig:portal}) offre une interface simplifiée aux
clients externes : consultation de l'avancement, acceptation de devis, signature
de contrats et messagerie avec l'équipe projet.

\medskip

L'interface portail est volontairement distincte de l'ERP interne : navigation
réduite à quatre entrées (tableau de bord, discussions, calendrier, devis), sans
accès aux modules de gestion internes (projets, facturation, configuration).
Les routes React sont regroupées sous le préfixe \texttt{/portal/*} et protégées
par le middleware d'authentification ; un utilisateur possédant le rôle
\texttt{client\_extern} est automatiquement redirigé vers le portail après
connexion, tandis qu'un collaborateur interne ne peut pas y accéder.

\medskip

\textbf{Tableau de bord client.} L'écran principal (Figure~\ref{fig:portal})
présente trois blocs d'information. Le widget \textit{Avancement du chantier}
affiche un indicateur circulaire (\texttt{ProgressRing}) alimenté par le champ
\texttt{progress\_percent} du projet, accompagné de la référence chantier, du
statut, de la ville et de la date de fin prévue. Le panneau \textit{Étapes de
la semaine} liste les jalons et tâches à échéance sur les sept prochains jours,
filtrés par projet sélectionné. Enfin, le fil \textit{Chantier en direct}
(\texttt{ChantierLiveFeed}) expose les photos terrain publiées par l'équipe
CIVCO BTP, rafraîchies périodiquement pour offrir une visibilité continue
sur l'avancement des travaux.

\medskip

\textbf{Parcours documentaire et contractuel.} Depuis l'onglet \textit{Devis},
le client consulte les propositions commerciales qui le concernent et peut
accepter un devis en un clic ; le statut est immédiatement mis à jour côté
serveur. Pour les contrats, le composant \texttt{ClientContractSignature}
intègre un pad de signature tactile (canvas HTML5) : la signature est capturée,
transmise à l'API et associée au document contractuel, déclenchant la mise à
jour du statut et la notification de l'équipe projet.

\medskip

\textbf{Provisionnement et sécurité.} L'accès portail est créé par
l'administrateur depuis la fiche client (module Clients, Chapitre~4) : un
identifiant de connexion est généré, associé au client concerné via le champ
\texttt{client\_id} de l'utilisateur, et activé ou désactivé à la demande.
Toutes les requêtes API du portail sont filtrées par ce lien client $\rightarrow$
projet : un client externe ne voit que ses propres chantiers, devis et
échanges, conformément au principe de moindre privilège défini au Chapitre~3.

\subsection{Configuration}

\begin{figure}[H]
  \centering
  \includegraphics[width=\textwidth]{config.png}
  \caption{Module Configuration : rôles, modèles de documents, contrôles d'impression}
  \label{fig:config}
\end{figure}

Le module Configuration (Figure~\ref{fig:config}) est réservé à l'administrateur :
gestion des rôles et permissions, modèles de contrats et documents, paramètres
d'impression, thème visuel et référentiels (badges, lots, secteurs).

\medskip

Le panneau de configuration est organisé en onglets thématiques : rôles et permissions,
modèles de contrats, contrôles d'impression, modèles de documents et paramètres
généraux. L'administrateur peut créer un nouveau rôle métier, lui attribuer un
sous-ensemble de permissions et l'assigner aux utilisateurs concernés, sans
intervention technique sur le code.

% ─────────────────────────────────────────────────────────────────────
\section{Tests et validation}

La validation de l'application s'appuie sur deux dispositifs complémentaires : une
\textbf{suite de tests automatisés} PHPUnit au niveau du backend, et une
\textbf{campagne de tests manuels} couvrant les parcours utilisateurs de bout en bout.

\subsection{Tests automatisés backend}

Le projet inclut une suite de tests d'intégration PHPUnit dans \texttt{tests/Feature/} :

\begin{description}
  \item[AuthTest] Connexion, déconnexion, accès au profil utilisateur, rejet des
  identifiants invalides.
  \item[ClientProjectTest] CRUD clients et projets, relations, contraintes d'accès.
  \item[QuoteInvoiceTest] Création de devis, lignes de détail, conversion en facture,
  calcul des totaux.
  \item[DocumentTest] Upload et gestion des documents attachés aux projets.
  \item[ExpenseTest] Enregistrement des dépenses de chantier.
  \item[DashboardTest] Agrégation des indicateurs du tableau de bord.
  \item[RoleAccessTest] Vérification des redirections post-login et des permissions
  par profil (comptable, commercial, chef de projet, administrateur).
\end{description}

Les tests utilisent \texttt{RefreshDatabase} et les seeders de démonstration pour
garantir un environnement isolé et reproductible. La commande \texttt{php artisan test}
exécute l'ensemble de la suite.

\subsection{Campagne de tests manuels}

Une checklist de tests manuels a été établie pour valider les parcours que les tests
automatisés ne couvrent pas entièrement :

\begin{itemize}
  \item \textbf{Authentification} : connexion, déconnexion, expiration de session,
  redirection selon le profil.
  \item \textbf{Projets} : création complète (client, phases, équipe, géolocalisation),
  modification, consultation par un utilisateur non-admin.
  \item \textbf{Documents commerciaux} : cycle devis $\rightarrow$ acceptation $\rightarrow$
  facture $\rightarrow$ paiement ; impression avec et sans en-tête.
  \item \textbf{Contrats} : compilation depuis un modèle, signature portail client.
  \item \textbf{Tâches} : création, assignation, changement de statut, visibilité
  selon les permissions.
  \item \textbf{Portail client} : connexion client, consultation projet, acceptation
  devis, messagerie.
  \item \textbf{Configuration} : création de rôle, attribution de permissions,
  modification du thème.
  \item \textbf{Responsive} : consultation sur poste de bureau et tablette.
\end{itemize}

\subsection{Validation qualitative}

L'application a été validée en conditions réelles au sein de CIVCO BTP, avec les
observations suivantes :

\begin{itemize}
  \item \textbf{Adéquation métier} : les modules couvrent les processus quotidiens
  identifiés au Chapitre~1 (suivi de chantier, production documentaire, planification).
  \item \textbf{Ergonomie} : l'interface sombre et la navigation par sidebar ont été
  appréciées par les testeurs internes pour un usage prolongé.
  \item \textbf{Performance} : les listes et fiches détail se chargent en moins de
  500~ms en environnement local ; le tableau de bord agrège les indicateurs sans
  latence perceptible.
  \item \textbf{Sécurité} : les tests de contournement de permissions (accès direct
  à une URL sans droit) sont correctement bloqués côté API avec une réponse
  \texttt{403 Forbidden}.
  \item \textbf{Impression} : le mécanisme de suivi des copies officielles et le
  filigrane « COPIE - NON OFFICIEL » fonctionnent conformément aux exigences métier.
\end{itemize}

\subsection{Synthèse des tests}

Le Tableau~\ref{tab:test_summary} récapitule la couverture de validation par module.

\begin{table}[H]
\caption{Synthèse de la couverture de tests par module}
\label{tab:test_summary}
\centering
\small
\begin{tabular}{|l|c|c|p{5.5cm}|}
\hline
\textbf{Module} & \textbf{Auto.} & \textbf{Manuel} & \textbf{Points vérifiés} \\
\hline
Authentification & $\checkmark$ & $\checkmark$ & Login, logout, session, redirection \\
\hline
Projets & $\checkmark$ & $\checkmark$ & CRUD, phases, équipe, carte \\
\hline
Clients & $\checkmark$ & $\checkmark$ & CRUD, archivage, portail \\
\hline
Devis / Factures & $\checkmark$ & $\checkmark$ & Lignes, conversion, impression \\
\hline
Tâches & --- & $\checkmark$ & Assignation, statuts, permissions \\
\hline
Contrats & --- & $\checkmark$ & Compilation, signature, avenants \\
\hline
Configuration & --- & $\checkmark$ & Rôles, modèles, thème \\
\hline
Portail client & --- & $\checkmark$ & Consultation, acceptation, messagerie \\
\hline
\end{tabular}
\end{table}

% ─────────────────────────────────────────────────────────────────────
\section{Difficultés rencontrées et solutions}

\subsection{Gestion des permissions granulaires}

La mise en place du \textbf{RBAC} avec plus de trente permissions a nécessité une
coordination entre le backend (middleware, seeders) et le frontend (PermissionGate,
routePermissions). La solution retenue centralise la résolution des permissions
dans \texttt{AuthContextService} côté API et \texttt{permissionResolver.js} côté
React, garantissant une source de vérité unique.

\subsection{Impression documentaire}

La génération de PDF côté client (composants React dédiés) et le suivi des quotas
d'impression officielle ont demandé un mécanisme de comptage persistant
(\texttt{PrintTrackingService}) et une interface de choix avant chaque impression
(\texttt{PrintOptionsModal}).

\subsection{Géolocalisation des chantiers}

L'intégration de la carte Leaflet avec géocodage automatique des adresses de chantier
a été réalisée via \texttt{SiteGeocodingService}, qui interroge l'API Nominatim avec
un repli sur un référentiel de coordonnées des villes marocaines
(\texttt{MoroccoCityCoordinates}) en cas d'échec.

\subsection{Synchronisation frontend / backend}

La cohérence entre les permissions affichées dans l'interface et celles vérifiées
côté API a nécessité un travail de synchronisation : chaque route frontend dans
\texttt{routePermissions.js} doit correspondre à un contrôle \texttt{permission:*}
sur l'endpoint associé. Les tests \texttt{RoleAccessTest} et les tests manuels
multi-profils ont permis de détecter et corriger les écarts.

\section*{Conclusion}

Ce chapitre a présenté la réalisation concrète de l'application de gestion de
projets BTP développée pour CIVCO BTP. Au-delà de la conception formalisée au
chapitre précédent, nous avons décrit l'environnement d'exécution, l'organisation
du code, les modules livrés et les dispositifs de validation retenus pour garantir
la fiabilité de la solution en conditions d'usage réelles.

\medskip

\textbf{Environnement et architecture.} L'application repose sur une stack
\textbf{Laravel 13 / React 19 / Vite / MySQL}, orchestrée en développement par
deux processus parallèles (\texttt{php artisan serve} et \texttt{npm run dev}).
Le proxy Vite assure la communication transparente entre la SPA et l'API, tandis
que \textbf{Laravel Sanctum} gère l'authentification par session sécurisée.
Cette séparation frontend / backend a permis de maintenir une frontière nette :
l'interface React ne contient aucune règle métier critique ; la validation, les
calculs financiers et la persistance sont systématiquement assurés côté serveur.

\medskip

\textbf{Implémentation backend.} L'API REST expose plus de quarante contrôleurs
sous le préfixe \texttt{/api/v1/}, structurés par domaine (projets, clients,
devis, factures, tâches, contrats, configuration). Chaque endpoint sensible est
protégé par un middleware de permission (\texttt{permission:*}) et validé par
un Form Request dédié. Les services métier (\texttt{PrintTrackingService},
\texttt{DocumentTemplateCompilationService}, \texttt{TeamMemberService}, etc.)
centralisent les règles de gestion, tandis que les API Resources normalisent
les réponses JSON consommées par le frontend.

\medskip

\textbf{Interface utilisateur.} Le frontend React est organisé en modules par
domaine fonctionnel, avec une sidebar filtrée selon le profil connecté. Les
écrans principaux ont été implémentés et illustrés dans ce chapitre :
authentification (Figure~\ref{fig:login}), tableau de bord
(Figure~\ref{fig:dashboard}), gestion des projets et des clients
(Figures~\ref{fig:projects} et~\ref{fig:clients}), documents commerciaux
(Figures~\ref{fig:quotes} et~\ref{fig:invoices}), planification des tâches
(Figure~\ref{fig:tasks}), carte géographique des chantiers (Figure~\ref{fig:map}),
portail client (Figure~\ref{fig:portal}) et module de configuration
(Figure~\ref{fig:config}). Cette couverture correspond au périmètre fonctionnel
défini au Chapitre~3 et aux processus quotidiens de CIVCO BTP.

\medskip

\textbf{Validation et déploiement.} La qualité de la solution est assurée par
une suite de tests PHPUnit (authentification, CRUD métier, conversion
devis $\rightarrow$ facture, contrôle d'accès par rôle) complétée par une
campagne de tests manuels sur les parcours de bout en bout. L'étude comparative
d'hébergement a orienté le choix vers \textbf{o2switch}, retenu pour sa simplicité
de mise en production, son support francophone et son rapport qualité/prix adapté
à une PME du secteur BTP.

\medskip

\textbf{Bilan.} Les principales difficultés --- permissions granulaires,
impression documentaire, géolocalisation des chantiers et synchronisation
frontend / backend --- ont été résolues par des services dédiés et une
source de vérité unique côté API. L'application est fonctionnelle, testée et
prête pour un déploiement chez CIVCO BTP. Le chapitre suivant conclut ce rapport
en synthétisant les apports du projet et en esquissant les perspectives
d'évolution (notifications temps réel, application mobile, intégration comptable).

\newpage
